%------------------------------------------------------------------------------%
\addtocounter{exemplo}{1}
\begin{frame}{Exemplo \theexemplo}
  Encontre os máximos e mínimos absolutos da função
  \[
    f(x, y) = x^2 + 2y^2 - x
  \]
  na região
  \[
    x^2 + y^2 \leq 1
  \]
\end{frame}

%------------------------------------------------------------------------------%
\begin{frame}{Exemplo \theexemplo}
  Pontos críticos interiores
  \[
    f_x(x, y)
    \pause
    = \D{}{x}\left(x^2 + 2y^2 - x\right)
    \pause
    = 2x-1
    \pause
    = 0
    \pause
    \quad\Rightarrow\quad
    x = \frac{1}{2}
  \]
  \[
    \pause
    f_y(x, y)
    \pause
    = \D{}{x}\left(x^2 + 2y^2 - x\right)
    \pause
    = 4y
    \pause
    = 0
    \pause
    \quad\Rightarrow\quad
    y = 0
  \]
  \pause
  É um ponto interior

  \pause
  \vspace{\baselineskip}
  Valor da função
  \[
    f\left(\sfrac{1}{2}, 0\right)
    \pause
    = \left(x^2 + 2y^2 - x\right)\evalat{\left(\sfrac{1}{2}, 0\right)}{}
    \pause
    = \left(\frac{1}{2}\right)^2 + 2\times 0 ^2 - \frac{1}{2}
    \pause
    = \frac{1}{4} - \frac{2}{4}
    \pause
    = -\frac{1}{4}
  \]
\end{frame}

%------------------------------------------------------------------------------%
\begin{frame}{Exemplo \theexemplo}
  Fronteira
  \[
    x^2+y^2 = 1
    \pause
    \quad\Rightarrow\quad
    y^2 = 1 - x^2
  \]
  \pause
  com \(-1\leq x \leq 1\)

  \pause
  \vspace{\baselineskip}
  Função $f$ restrita à fronteira
  \[
    g(x)
    \pause
    = x^2 + 2\left(1 - x^2\right) -x
    \pause
    = x^2 + 2 - 2x^2 - x
    \pause
    = 2 - x^2 - x
  \]
\end{frame}

%------------------------------------------------------------------------------%
\begin{frame}{Exemplo \theexemplo}
  Candidatos a máximos e mínimos na fronteira

  \pause
  \vspace{\baselineskip}
  Extremidade esquerda
  \[
    x = -1
    \pause
    \quad\Rightarrow\quad
    \pause
    y^2 = 1 - (-1)^2 = 0
  \]
  \pause
  Extremidade direita
  \[
    \pause
    x = 1
    \pause
    \quad\Rightarrow\quad
    y^2 = 1 - 1^2
    \pause
    = 0
  \]
\end{frame}

%------------------------------------------------------------------------------%
\begin{frame}{Exemplo \theexemplo}
  Interior
  \[
    \pause
    g'(x)
    \pause
    = \frac{d}{dx}\left(2 - x^2 - x\right)
    \pause
    = -2x - 1
    \pause
    = 0
    \pause
    \quad\Rightarrow\quad
    x = -\frac{1}{2}
  \]
  \[
    \pause
    y^2 = 1 - x^2
    \pause
    = 1 - \left(-\frac{1}{2}\right)^2
    \pause
    = \frac{4}{4} - \frac{1}{4}
    \pause
    = \frac{3}{4}
 \]
 \[
    \pause
    \abs{y} = \sqrt{\frac{3}{4}}
  \]
  \[
    \pause
    y
    = \pm\sqrt{\frac{3}{4}}
    \pause
    = \pm\frac{\sqrt{3}}{2}
  \]
\end{frame}

%------------------------------------------------------------------------------%
\begin{frame}{Exemplo \theexemplo}
  Avaliando a função nos pontos
  \[
    \pause
    f(-1, 0)
    \pause
    = \left(x^2 + 2y^2 - x\right)\evalat{(-1, 0)}{}
    \pause
    = 1 + 0 - (-1)
    \pause
    = 2
  \]
  \[
    \pause
    f(1, 0)
    \pause
    = \left(x^2 + 2y^2 - x\right)\evalat{(1, 0)}{}
    \pause
    = 1 + 0 - 1
    \pause
    = 0
  \]
\end{frame}

%------------------------------------------------------------------------------%
\begin{frame}{Exemplo \theexemplo}
  \[
    f\left(-\frac{1}{2}, \frac{\sqrt{3}}{2}\right)
    \pause
    = \left(x^2 + 2y^2 - x\right)\evalat{\left(-\frac{1}{2}, \frac{\sqrt{3}}{2}\right)}{}
    \pause
    = \left(-\frac{1}{2}\right)^2 + 2\left(\frac{\sqrt{3}}{2}\right)^2 - \left(-\frac{1}{2}\right)
  \]
  \[
    \pause
    \hphantom{f\left(-\frac{1}{2}, \frac{\sqrt{3}}{2}\right)}\,
    = \frac{1}{4} + 2\frac{3}{4} + \frac{2}{4}
    \pause
    = \frac{9}{4}
    = 2,25
  \]

  \pause
  \[
    f\left(-\frac{1}{2}, -\frac{\sqrt{3}}{2}\right)
    \pause
    = \left(x^2 + 2y^2 - x\right)\evalat{\left(-\frac{1}{2}, -\frac{\sqrt{3}}{2}\right)}{}
    \pause
    = \left(-\frac{1}{2}\right)^2 + 2\left(-\frac{\sqrt{3}}{2}\right)^2 - \left(-\frac{1}{2}\right)
  \]
  \[
    \pause
    \hphantom{f\left(-\frac{1}{2}, -\frac{\sqrt{3}}{2}\right)}\,
    = \frac{1}{4} + 2\frac{3}{4} + \frac{2}{4}
    \pause
    = \frac{9}{4}
    = 2,25
  \]
\end{frame}

%------------------------------------------------------------------------------%
\begin{frame}{Exemplo \theexemplo}
  Máximo absoluto de $f$ na região é \(\frac{9}{4}\) e ocorre nos pontos \\
  \(\left(-\frac{1}{2}, -\frac{\sqrt{3}}{2}\right)\) e
  \(\left(-\frac{1}{2},  \frac{\sqrt{3}}{2}\right)\)

  \vspace{\baselineskip}
  Mínimo absoluto de $f$ na região é \(-\frac{1}{4}\)
  e ocorre no ponto
  \(\left(\frac{1}{2}, 0\right)\)
\end{frame}

%------------------------------------------------------------------------------%
